\documentclass{amsart}
% For dealing with references we use the comment environment
\usepackage{verbatim}
\newenvironment{reference}{\comment}{\endcomment}
%\newenvironment{reference}{}{}
\newenvironment{slogan}{\comment}{\endcomment}
\newenvironment{history}{\comment}{\endcomment}

% For commutative diagrams we use Xy-pic
\usepackage[all]{xy}

% We use 2cell for 2-commutative diagrams.
\xyoption{2cell}
\UseAllTwocells

% We use multicol for the list of chapters between chapters
\usepackage{multicol}

% This is generally recommended for better output
\usepackage{lmodern}
\usepackage[T1]{fontenc}

% For cross-file-references

% Package for hypertext links:
\usepackage{hyperref}

% For any local file, say "hello.tex" you want to link to please

% Theorem environments.
%
\theoremstyle{plain}
\newtheorem{theorem}[subsection]{Theorem}
\newtheorem{proposition}[subsection]{Proposition}
\newtheorem{lemma}[subsection]{Lemma}

\theoremstyle{definition}
\newtheorem{definition}[subsection]{Definition}
\newtheorem{example}[subsection]{Example}
\newtheorem{exercise}[subsection]{Exercise}
\newtheorem{situation}[subsection]{Situation}

\theoremstyle{remark}
\newtheorem{remark}[subsection]{Remark}
\newtheorem{remarks}[subsection]{Remarks}

\numberwithin{equation}{subsection}

% Macros
%
\def\lim{\mathop{\mathrm{lim}}\nolimits}
\def\colim{\mathop{\mathrm{colim}}\nolimits}
\def\Spec{\mathop{\mathrm{Spec}}}
\def\Hom{\mathop{\mathrm{Hom}}\nolimits}
\def\Ext{\mathop{\mathrm{Ext}}\nolimits}
\def\SheafHom{\mathop{\mathcal{H}\!\mathit{om}}\nolimits}
\def\SheafExt{\mathop{\mathcal{E}\!\mathit{xt}}\nolimits}
\def\Sch{\mathit{Sch}}
\def\Mor{\mathop{\mathrm{Mor}}\nolimits}
\def\Ob{\mathop{\mathrm{Ob}}\nolimits}
\def\Sh{\mathop{\mathit{Sh}}\nolimits}
\def\NL{\mathop{N\!L}\nolimits}
\def\CH{\mathop{\mathrm{CH}}\nolimits}
\def\proetale{{pro\text{-}\acute{e}tale}}
\def\etale{{\acute{e}tale}}
\def\QCoh{\mathit{QCoh}}
\def\Ker{\mathop{\mathrm{Ker}}}
\def\Im{\mathop{\mathrm{Im}}}
\def\Coker{\mathop{\mathrm{Coker}}}
\def\Coim{\mathop{\mathrm{Coim}}}

% Boxtimes
%
\DeclareMathSymbol{\boxtimes}{\mathbin}{AMSa}{"02}

%
% Macros for moduli stacks/spaces
%
\def\QCohstack{\mathcal{QC}\!\mathit{oh}}
\def\Cohstack{\mathcal{C}\!\mathit{oh}}
\def\Spacesstack{\mathcal{S}\!\mathit{paces}}
\def\Quotfunctor{\mathrm{Quot}}
\def\Hilbfunctor{\mathrm{Hilb}}
\def\Curvesstack{\mathcal{C}\!\mathit{urves}}
\def\Polarizedstack{\mathcal{P}\!\mathit{olarized}}
\def\Complexesstack{\mathcal{C}\!\mathit{omplexes}}
% \Pic is the operator that assigns to X its picard group, usage \Pic(X)
% \Picardstack_{X/B} denotes the Picard stack of X over B
% \Picardfunctor_{X/B} denotes the Picard functor of X over B
\def\Pic{\mathop{\mathrm{Pic}}\nolimits}
\def\Picardstack{\mathcal{P}\!\mathit{ic}}
\def\Picardfunctor{\mathrm{Pic}}
\def\Deformationcategory{\mathcal{D}\!\mathit{ef}}

\setlength{\emergencystretch}{3em}
\begin{document}
\title[Stacks Project, Tag 0D1W]{276 --- Every etale structure-sheaf module on a zero-dimensional local scheme is quasi-coherent}
\maketitle
\noindent Copyright (C) 2005--2025 Johan de Jong. Stacks Project authors. Source: \href{https://stacks.math.columbia.edu/tag/0D1W}{Tag 0D1W}.

\noindent Editorial difficulty note (not author text). Difficulty H1: An arbitrary module sheaf must be reconstructed from global sections, not merely checked on a given affine presentation. The proof reduces every etale object to finite local etale algebras and uses semilinear Galois descent twice to handle both Galois and non-Galois residue extensions..

\noindent Editorial provenance note (not author text). History: excerpt from revision \texttt{a0444\allowbreak 6e57e\allowbreak c1fbc\allowbreak 252a8\allowbreak 71afc\allowbreak ec775\allowbreak 2fb28\allowbreak 07b14}, etale-cohomology.tex, lines 9176--9238. Reference presentation was adapted; original-author context and core support blocks were retained or added. The mathematical wording is unchanged. Licensed under GNU FDL 1.2 or later; no invariant sections or cover texts. See ../licenses/COPYING. Context blocks reproduce author definitions and supporting results; they are not additional counted problems.

\begin{lemma}
\label{lemma-all-modules-quasi-coherent}
Let $R$ be a local ring of dimension $0$. Let $S = \Spec(R)$.
Then every $\mathcal{O}_S$-module on $S_\etale$ is quasi-coherent.
\end{lemma}

\begin{proof}
Let $\mathcal{F}$ be an $\mathcal{O}_S$-module on $S_\etale$.
We have to show that $\mathcal{F}$ is determined by the
$R$-module $M = \Gamma(S, \mathcal{F})$.
More precisely, if $\pi : X \to S$ is \'etale we have to
show that $\Gamma(X, \mathcal{F}) = \Gamma(X, \pi^*\widetilde{M})$.

\medskip\noindent
Let $\mathfrak m \subset R$ be the maximal ideal and let
$\kappa$ be the residue field. By
Algebra, Lemma \href{https://stacks.math.columbia.edu/tag/06RS}{lemma local dimension zero henselian (Tag 06RS)}
the local ring $R$ is henselian. If $X \to S$ is \'etale,
then the underlying topological space of $X$ is discrete
by Morphisms, Lemma \href{https://stacks.math.columbia.edu/tag/02GL}{lemma etale over field (Tag 02GL)}
and hence $X$ is a disjoint union of affine schemes
each having one point. Moreover, if $X = \Spec(A)$ is affine and
has one point, then $R \to A$ is finite \'etale by
Algebra, Lemma \href{https://stacks.math.columbia.edu/tag/04GJ}{lemma mop up (Tag 04GJ)}.
We have to show that $\Gamma(X, \mathcal{F}) = M \otimes_R A$
in this case.

\medskip\noindent
The functor $A \mapsto A/\mathfrak m A$ defines an equivalence of
the category of finite \'etale $R$-algebras
with the category of finite separable $\kappa$-algebras by
Algebra, Lemma \href{https://stacks.math.columbia.edu/tag/04GK}{lemma henselian cat finite etale (Tag 04GK)}.
Let us first consider the case where $A/\mathfrak m A$
is a Galois extension of $\kappa$ with Galois group $G$.
For each $\sigma \in G$ let $\sigma : A \to A$ denote the
corresponding automorphism of $A$ over $R$.
Let $N = \Gamma(X, \mathcal{F})$.
Then $\Spec(\sigma) : X \to X$ is an automorphism over $S$
and hence pullback by this defines a map $\sigma : N \to N$
which is a $\sigma$-linear map: $\sigma(an) = \sigma(a) \sigma(n)$
for $a \in A$ and $n \in N$.
We will apply Galois descent to the quasi-coherent module
$\widetilde{N}$ on $X$ endowed with the isomorphisms
coming from the action on $\sigma$ on $N$. See Descent, Lemma
\href{https://stacks.math.columbia.edu/tag/0D1V}{lemma galois descent more general (Tag 0D1V)}.
This lemma tells us there is an isomorphism $N = N^G \otimes_R A$.
On the other hand, it is clear that $N^G = M$ by the sheaf property
for $\mathcal{F}$. Thus the required isomorphism holds.

\medskip\noindent
The general case (with $A$ local and finite \'etale over $R$)
is deduced from the Galois case as follows. Choose $A \to B$
finite \'etale such that $B$ is local with residue field
Galois over $\kappa$. Let $G = \text{Aut}(B/R) = \text{Gal}(\kappa_B/\kappa)$.
Let $H \subset G$ be the Galois group corresponding to the
Galois extension $\kappa_B/\kappa_A$. Then as above one
shows that $\Gamma(X, \mathcal{F}) = \Gamma(\Spec(B), \mathcal{F})^H$.
By the result for Galois extensions (used twice) we get
$$
\Gamma(X, \mathcal{F}) = (M \otimes_R B)^H = M \otimes_R A
$$
as desired.
\end{proof}
\end{document}
